% Options for packages loaded elsewhere
\PassOptionsToPackage{unicode}{hyperref}
\PassOptionsToPackage{hyphens}{url}
\documentclass[
]{article}
\usepackage{xcolor}
\usepackage[margin=1in]{geometry}
\usepackage{amsmath,amssymb}
\setcounter{secnumdepth}{-\maxdimen} % remove section numbering
\usepackage{iftex}
\ifPDFTeX
  \usepackage[T1]{fontenc}
  \usepackage[utf8]{inputenc}
  \usepackage{textcomp} % provide euro and other symbols
\else % if luatex or xetex
  \usepackage{unicode-math} % this also loads fontspec
  \defaultfontfeatures{Scale=MatchLowercase}
  \defaultfontfeatures[\rmfamily]{Ligatures=TeX,Scale=1}
\fi
\usepackage{lmodern}
\ifPDFTeX\else
  % xetex/luatex font selection
\fi
% Use upquote if available, for straight quotes in verbatim environments
\IfFileExists{upquote.sty}{\usepackage{upquote}}{}
\IfFileExists{microtype.sty}{% use microtype if available
  \usepackage[]{microtype}
  \UseMicrotypeSet[protrusion]{basicmath} % disable protrusion for tt fonts
}{}
\makeatletter
\@ifundefined{KOMAClassName}{% if non-KOMA class
  \IfFileExists{parskip.sty}{%
    \usepackage{parskip}
  }{% else
    \setlength{\parindent}{0pt}
    \setlength{\parskip}{6pt plus 2pt minus 1pt}}
}{% if KOMA class
  \KOMAoptions{parskip=half}}
\makeatother
\usepackage{color}
\usepackage{fancyvrb}
\newcommand{\VerbBar}{|}
\newcommand{\VERB}{\Verb[commandchars=\\\{\}]}
\DefineVerbatimEnvironment{Highlighting}{Verbatim}{commandchars=\\\{\}}
% Add ',fontsize=\small' for more characters per line
\usepackage{framed}
\definecolor{shadecolor}{RGB}{248,248,248}
\newenvironment{Shaded}{\begin{snugshade}}{\end{snugshade}}
\newcommand{\AlertTok}[1]{\textcolor[rgb]{0.94,0.16,0.16}{#1}}
\newcommand{\AnnotationTok}[1]{\textcolor[rgb]{0.56,0.35,0.01}{\textbf{\textit{#1}}}}
\newcommand{\AttributeTok}[1]{\textcolor[rgb]{0.13,0.29,0.53}{#1}}
\newcommand{\BaseNTok}[1]{\textcolor[rgb]{0.00,0.00,0.81}{#1}}
\newcommand{\BuiltInTok}[1]{#1}
\newcommand{\CharTok}[1]{\textcolor[rgb]{0.31,0.60,0.02}{#1}}
\newcommand{\CommentTok}[1]{\textcolor[rgb]{0.56,0.35,0.01}{\textit{#1}}}
\newcommand{\CommentVarTok}[1]{\textcolor[rgb]{0.56,0.35,0.01}{\textbf{\textit{#1}}}}
\newcommand{\ConstantTok}[1]{\textcolor[rgb]{0.56,0.35,0.01}{#1}}
\newcommand{\ControlFlowTok}[1]{\textcolor[rgb]{0.13,0.29,0.53}{\textbf{#1}}}
\newcommand{\DataTypeTok}[1]{\textcolor[rgb]{0.13,0.29,0.53}{#1}}
\newcommand{\DecValTok}[1]{\textcolor[rgb]{0.00,0.00,0.81}{#1}}
\newcommand{\DocumentationTok}[1]{\textcolor[rgb]{0.56,0.35,0.01}{\textbf{\textit{#1}}}}
\newcommand{\ErrorTok}[1]{\textcolor[rgb]{0.64,0.00,0.00}{\textbf{#1}}}
\newcommand{\ExtensionTok}[1]{#1}
\newcommand{\FloatTok}[1]{\textcolor[rgb]{0.00,0.00,0.81}{#1}}
\newcommand{\FunctionTok}[1]{\textcolor[rgb]{0.13,0.29,0.53}{\textbf{#1}}}
\newcommand{\ImportTok}[1]{#1}
\newcommand{\InformationTok}[1]{\textcolor[rgb]{0.56,0.35,0.01}{\textbf{\textit{#1}}}}
\newcommand{\KeywordTok}[1]{\textcolor[rgb]{0.13,0.29,0.53}{\textbf{#1}}}
\newcommand{\NormalTok}[1]{#1}
\newcommand{\OperatorTok}[1]{\textcolor[rgb]{0.81,0.36,0.00}{\textbf{#1}}}
\newcommand{\OtherTok}[1]{\textcolor[rgb]{0.56,0.35,0.01}{#1}}
\newcommand{\PreprocessorTok}[1]{\textcolor[rgb]{0.56,0.35,0.01}{\textit{#1}}}
\newcommand{\RegionMarkerTok}[1]{#1}
\newcommand{\SpecialCharTok}[1]{\textcolor[rgb]{0.81,0.36,0.00}{\textbf{#1}}}
\newcommand{\SpecialStringTok}[1]{\textcolor[rgb]{0.31,0.60,0.02}{#1}}
\newcommand{\StringTok}[1]{\textcolor[rgb]{0.31,0.60,0.02}{#1}}
\newcommand{\VariableTok}[1]{\textcolor[rgb]{0.00,0.00,0.00}{#1}}
\newcommand{\VerbatimStringTok}[1]{\textcolor[rgb]{0.31,0.60,0.02}{#1}}
\newcommand{\WarningTok}[1]{\textcolor[rgb]{0.56,0.35,0.01}{\textbf{\textit{#1}}}}
\usepackage{graphicx}
\makeatletter
\newsavebox\pandoc@box
\newcommand*\pandocbounded[1]{% scales image to fit in text height/width
  \sbox\pandoc@box{#1}%
  \Gscale@div\@tempa{\textheight}{\dimexpr\ht\pandoc@box+\dp\pandoc@box\relax}%
  \Gscale@div\@tempb{\linewidth}{\wd\pandoc@box}%
  \ifdim\@tempb\p@<\@tempa\p@\let\@tempa\@tempb\fi% select the smaller of both
  \ifdim\@tempa\p@<\p@\scalebox{\@tempa}{\usebox\pandoc@box}%
  \else\usebox{\pandoc@box}%
  \fi%
}
% Set default figure placement to htbp
\def\fps@figure{htbp}
\makeatother
\setlength{\emergencystretch}{3em} % prevent overfull lines
\providecommand{\tightlist}{%
  \setlength{\itemsep}{0pt}\setlength{\parskip}{0pt}}
\usepackage{bookmark}
\IfFileExists{xurl.sty}{\usepackage{xurl}}{} % add URL line breaks if available
\urlstyle{same}
\hypersetup{
  pdftitle={TraCESahulMisc: Workflow Vignette},
  hidelinks,
  pdfcreator={LaTeX via pandoc}}

\title{TraCESahulMisc: Workflow Vignette}
\author{}
\date{\vspace{-2.5em}}

\begin{document}
\maketitle

\section{Introduction}\label{introduction}

This vignette provides a worked example using \texttt{TraCESahulMisc} to
load downscaled TraCE--Sahul climate data, summarise it in time,
calculate BIOCLIM variables, match palaeo-observations, generate
background points, and fit a quick Random Forest model to examine
habitat suitability through time.

\section{Installing and loading
packages}\label{installing-and-loading-packages}

For the vignette we need three packages.

\begin{Shaded}
\begin{Highlighting}[]
\FunctionTok{library}\NormalTok{(TraCESahulMisc)}
\FunctionTok{library}\NormalTok{(terra)}
\CommentTok{\#\textgreater{} terra 1.8.80}
\FunctionTok{library}\NormalTok{(randomForest)}
\CommentTok{\#\textgreater{} randomForest 4.7{-}1.2}
\CommentTok{\#\textgreater{} Type rfNews() to see new features/changes/bug fixes.}

\CommentTok{\# only for plotting}
\FunctionTok{library}\NormalTok{(rnaturalearth)}
\NormalTok{conts }\OtherTok{\textless{}{-}} \FunctionTok{vect}\NormalTok{(rnaturalearth}\SpecialCharTok{::}\FunctionTok{ne\_coastline}\NormalTok{(}\AttributeTok{scale =} \DecValTok{50}\NormalTok{))}
\NormalTok{conts }\OtherTok{\textless{}{-}} \FunctionTok{crop}\NormalTok{(conts, }\FunctionTok{ext}\NormalTok{(}\DecValTok{125}\NormalTok{, }\DecValTok{135}\NormalTok{, }\SpecialCharTok{{-}}\DecValTok{13}\NormalTok{, }\SpecialCharTok{{-}}\DecValTok{8}\NormalTok{))}
\end{Highlighting}
\end{Shaded}

\section{Downloading example data}\label{downloading-example-data}

\begin{Shaded}
\begin{Highlighting}[]
\NormalTok{base\_dir }\OtherTok{\textless{}{-}} \StringTok{"D:/TraCE{-}Sahul"}
\CommentTok{\# download\_trace\_data(base\_dir) \# Run once only}
\end{Highlighting}
\end{Shaded}

\section{Importing TraCE--Sahul data}\label{importing-tracesahul-data}

This small function will read in the downloaded data to an object called
\texttt{sahul} which will then have names for each of the variables.

\begin{Shaded}
\begin{Highlighting}[]
\NormalTok{vars }\OtherTok{\textless{}{-}} \FunctionTok{c}\NormalTok{(}\StringTok{"pr"}\NormalTok{, }\StringTok{"tasmax"}\NormalTok{, }\StringTok{"tasmin"}\NormalTok{)}

\NormalTok{load\_tracesahul }\OtherTok{\textless{}{-}} \ControlFlowTok{function}\NormalTok{(base, var) \{}
\NormalTok{  f1 }\OtherTok{\textless{}{-}} \FunctionTok{file.path}\NormalTok{(base, var,}
                  \FunctionTok{sprintf}\NormalTok{(}\StringTok{"TraCE\_22ka\_downscaled\_\%s\_decadal\_21k\_1500CE\_biascorr.nc"}\NormalTok{, var))}
\NormalTok{  f2 }\OtherTok{\textless{}{-}} \FunctionTok{file.path}\NormalTok{(base, var,}
                  \FunctionTok{sprintf}\NormalTok{(}\StringTok{"TraCE\_22ka\_downscaled\_\%s\_1500\_1990\_biascorr.nc"}\NormalTok{, var))}
  \FunctionTok{import\_TraCESahul}\NormalTok{(}\FunctionTok{c}\NormalTok{(f1, f2))}
\NormalTok{\}}
\NormalTok{sahul }\OtherTok{\textless{}{-}} \FunctionTok{lapply}\NormalTok{(vars, }\ControlFlowTok{function}\NormalTok{(v) }\FunctionTok{load\_tracesahul}\NormalTok{(base\_dir, v))}
\FunctionTok{names}\NormalTok{(sahul) }\OtherTok{\textless{}{-}}\NormalTok{ vars}
\NormalTok{sahul}
\CommentTok{\#\textgreater{} $pr}
\CommentTok{\#\textgreater{} class       : SpatRaster }
\CommentTok{\#\textgreater{} size        : 100, 200, 31740  (nrow, ncol, nlyr)}
\CommentTok{\#\textgreater{} resolution  : 0.05, 0.05  (x, y)}
\CommentTok{\#\textgreater{} extent      : 125, 135, {-}13, {-}8  (xmin, xmax, ymin, ymax)}
\CommentTok{\#\textgreater{} coord. ref. : lon/lat WGS 84 (CRS84) (OGC:CRS84) }
\CommentTok{\#\textgreater{} sources     : TraCE\_22ka\_downscaled\_pr\_decadal\_21k\_1500CE\_biascorr.nc  (25860 layers) }
\CommentTok{\#\textgreater{}               TraCE\_22ka\_downscaled\_pr\_1500\_1990\_biascorr.nc  (5880 layers) }
\CommentTok{\#\textgreater{} varnames    : pr (precipitation) }
\CommentTok{\#\textgreater{}               pr (precipitation) }
\CommentTok{\#\textgreater{} names       :       pr\_1,       pr\_2,       pr\_3,       pr\_4,       pr\_5,       pr\_6,        ... }
\CommentTok{\#\textgreater{} unit        : kg m{-}2 s{-}1 }
\CommentTok{\#\textgreater{} depth       : 1 to 12 (calendar month [calendar month]: 12 steps) }
\CommentTok{\#\textgreater{} time (raw)  : {-}20045 to 1989 (2645 steps) }
\CommentTok{\#\textgreater{} }
\CommentTok{\#\textgreater{} $tasmax}
\CommentTok{\#\textgreater{} class       : SpatRaster }
\CommentTok{\#\textgreater{} size        : 100, 200, 31740  (nrow, ncol, nlyr)}
\CommentTok{\#\textgreater{} resolution  : 0.05, 0.05  (x, y)}
\CommentTok{\#\textgreater{} extent      : 125, 135, {-}13, {-}8  (xmin, xmax, ymin, ymax)}
\CommentTok{\#\textgreater{} coord. ref. : lon/lat WGS 84 (CRS84) (OGC:CRS84) }
\CommentTok{\#\textgreater{} sources     : TraCE\_22ka\_downscaled\_tasmax\_decadal\_21k\_1500CE\_biascorr.nc  (25860 layers) }
\CommentTok{\#\textgreater{}               TraCE\_22ka\_downscaled\_tasmax\_1500\_1990\_biascorr.nc  (5880 layers) }
\CommentTok{\#\textgreater{} varnames    : tasmax (Daily Maximum Near{-}Surface Air Temperatures) }
\CommentTok{\#\textgreater{}               tasmax (Daily Maximum Near{-}Surface Air Temperatures) }
\CommentTok{\#\textgreater{} names       : tasmax\_1, tasmax\_2, tasmax\_3, tasmax\_4, tasmax\_5, tasmax\_6,     ... }
\CommentTok{\#\textgreater{} unit        : Celcius }
\CommentTok{\#\textgreater{} depth       : 1 to 12 (calendar month [calendar month]: 12 steps) }
\CommentTok{\#\textgreater{} time (raw)  : {-}20045 to 1989 (2645 steps) }
\CommentTok{\#\textgreater{} }
\CommentTok{\#\textgreater{} $tasmin}
\CommentTok{\#\textgreater{} class       : SpatRaster }
\CommentTok{\#\textgreater{} size        : 100, 200, 31740  (nrow, ncol, nlyr)}
\CommentTok{\#\textgreater{} resolution  : 0.05, 0.05  (x, y)}
\CommentTok{\#\textgreater{} extent      : 125, 135, {-}13, {-}8  (xmin, xmax, ymin, ymax)}
\CommentTok{\#\textgreater{} coord. ref. : lon/lat WGS 84 (CRS84) (OGC:CRS84) }
\CommentTok{\#\textgreater{} sources     : TraCE\_22ka\_downscaled\_tasmin\_decadal\_21k\_1500CE\_biascorr.nc  (25860 layers) }
\CommentTok{\#\textgreater{}               TraCE\_22ka\_downscaled\_tasmin\_1500\_1990\_biascorr.nc  (5880 layers) }
\CommentTok{\#\textgreater{} varnames    : tasmin (Daily Minimum Near{-}Surface Air Temperatures) }
\CommentTok{\#\textgreater{}               tasmin (Daily Minimum Near{-}Surface Air Temperatures) }
\CommentTok{\#\textgreater{} names       : tasmin\_1, tasmin\_2, tasmin\_3, tasmin\_4, tasmin\_5, tasmin\_6,     ... }
\CommentTok{\#\textgreater{} unit        : Celcius }
\CommentTok{\#\textgreater{} depth       : 1 to 12 (calendar month [calendar month]: 12 steps) }
\CommentTok{\#\textgreater{} time (raw)  : {-}20045 to 1989 (2645 steps)}
\end{Highlighting}
\end{Shaded}

\section{Temporal summarisation}\label{temporal-summarisation}

Here we can use the \texttt{summarise\_TraCESahul} to aggregate the data
in time.

The function is able to summarise an imported TraCESahul SpatRaster to
annual, monthly, or seasonal climatologies.

Monthly and seasonal summaries can be taken across the full record or
grouped into right-aligned windows of years. When
\texttt{window\ =\ 10}, the pre-1500 data are already decadal
climatologies and so are treated accordingly when deriving annual,
monthly, or seasonal outputs. The windows are guaranteed not to mix data
from before and after 1500 CE as there is a switch from decadal monthly
climatologies to annual monthly data from 1500 CE onwards.

Seasonal summaries follow austral seasons. December is treated as
belonging to the following year/timestep to ensure correct DJF
alignment, including when time steps are spaced irregularly or when
windowed summaries are requested.

Below, the imported data is aggregated to monthly averages using a
30-year right aligned window.

\begin{Shaded}
\begin{Highlighting}[]
\NormalTok{pr\_monthly\_win }\OtherTok{\textless{}{-}} \FunctionTok{summarise\_TraCESahul}\NormalTok{(sahul}\SpecialCharTok{$}\NormalTok{pr, }\AttributeTok{type =} \StringTok{"monthly"}\NormalTok{,}
                                       \AttributeTok{sumfun =} \StringTok{"mean"}\NormalTok{, }\AttributeTok{window =} \DecValTok{30}\NormalTok{)}
\NormalTok{tasmax\_monthly\_win }\OtherTok{\textless{}{-}} \FunctionTok{summarise\_TraCESahul}\NormalTok{(sahul}\SpecialCharTok{$}\NormalTok{tasmax, }\AttributeTok{type =} \StringTok{"monthly"}\NormalTok{,}
                                           \AttributeTok{sumfun =} \StringTok{"mean"}\NormalTok{, }\AttributeTok{window =} \DecValTok{30}\NormalTok{)}
\NormalTok{tasmin\_monthly\_win }\OtherTok{\textless{}{-}} \FunctionTok{summarise\_TraCESahul}\NormalTok{(sahul}\SpecialCharTok{$}\NormalTok{tasmin, }\AttributeTok{type =} \StringTok{"monthly"}\NormalTok{,}
                                           \AttributeTok{sumfun =} \StringTok{"mean"}\NormalTok{, }\AttributeTok{window =} \DecValTok{30}\NormalTok{)}

\NormalTok{lyr\_names }\OtherTok{\textless{}{-}} \FunctionTok{rep}\NormalTok{(month.abb, }\FunctionTok{nlyr}\NormalTok{(pr\_monthly\_win)}\SpecialCharTok{/}\DecValTok{12}\NormalTok{)}
\FunctionTok{names}\NormalTok{(pr\_monthly\_win) }\OtherTok{\textless{}{-}} \FunctionTok{paste0}\NormalTok{(lyr\_names, }\StringTok{"\_pr"}\NormalTok{)}
\FunctionTok{names}\NormalTok{(tasmax\_monthly\_win) }\OtherTok{\textless{}{-}} \FunctionTok{paste0}\NormalTok{(lyr\_names, }\StringTok{"\_tasmax"}\NormalTok{)}
\FunctionTok{names}\NormalTok{(tasmin\_monthly\_win) }\OtherTok{\textless{}{-}} \FunctionTok{paste0}\NormalTok{(lyr\_names, }\StringTok{"\_tasmin"}\NormalTok{)}
\end{Highlighting}
\end{Shaded}

\section{BIOCLIM calculation}\label{bioclim-calculation}

The \texttt{bioclim\_TraCESahul} function calculates bioclimatic
variables (BIO1--BIO19) from monthly TraCESahul climate data using
\texttt{fastbioclim::derive\_bioclim}.

The function requires monthly maximum temperature, minimum temperature,
and precipitation rasters, performs optional precipitation scaling,
processes each year independently, and returns either a list of annual
bioclim stacks or collated multi-year rasters for each requested
bioclimatic variable.

Ideally, the data should be the output from
\texttt{summarise\_TraCESahul} as there are a number of checks performed
on the data, which are guaranteed to pass if the data has been
aggregated with \texttt{summarise\_TraCESahul}.

Precipitation data is automatically converted from a flux (i.e.~kg m-2
s-1) to a rate (mm/day) if the units are set. See the \texttt{pr\_scale}
and \texttt{new\_units} arguments for further details.

Data is automatically saved into a output directory that can be
specified. If not specified a \texttt{tempdir()} is created and the user
is notified of the location of the directory so files can be retrieved.

Optionally, the data can be collated so that each item in the list that
is returned is one variable. The alternative \texttt{collate\ =\ FALSE}
results in each item in the list being all the bioclim variables for a
given year.

\begin{Shaded}
\begin{Highlighting}[]
\NormalTok{bioclims }\OtherTok{\textless{}{-}} \FunctionTok{bioclim\_TraCESahul}\NormalTok{(}
  \AttributeTok{tasmax =}\NormalTok{ tasmax\_monthly\_win,}
  \AttributeTok{tasmin =}\NormalTok{ tasmin\_monthly\_win,}
  \AttributeTok{pr =}\NormalTok{ pr\_monthly\_win,}
  \CommentTok{\# store in a folder in basedir. Not recommended}
  \AttributeTok{outdir =} \FunctionTok{paste0}\NormalTok{(base\_dir,}\StringTok{"/bioclims"}\NormalTok{),}
  \AttributeTok{bioclims =} \FunctionTok{c}\NormalTok{(}\DecValTok{1}\NormalTok{, }\DecValTok{4}\NormalTok{, }\DecValTok{5}\NormalTok{, }\DecValTok{6}\NormalTok{, }\DecValTok{7}\NormalTok{, }\DecValTok{12}\NormalTok{, }\DecValTok{13}\NormalTok{, }\DecValTok{14}\NormalTok{, }\DecValTok{17}\NormalTok{, }\DecValTok{18}\NormalTok{), }\CommentTok{\# default is to use 1:19}
  \AttributeTok{collate =} \ConstantTok{TRUE}\NormalTok{)}
\NormalTok{bioclims}
\end{Highlighting}
\end{Shaded}

A quick plot shows the first item in the list is all the BIO01 (mean
annual temperature) results.

\begin{center}\includegraphics[width=1\linewidth]{TraCESahulMisc_workflow_files/figure-latex/plotBio-1} \end{center}

\section{Generate a mask}\label{generate-a-mask}

Here we generate the mask required by the fossil matching function.

\begin{Shaded}
\begin{Highlighting}[]
\CommentTok{\# make a mask using bio01}
\NormalTok{mask\_ras }\OtherTok{\textless{}{-}}\NormalTok{ bioclims}\SpecialCharTok{$}\NormalTok{bio01}
\NormalTok{mask\_ras }\OtherTok{\textless{}{-}} \FunctionTok{ifel}\NormalTok{(}\FunctionTok{is.na}\NormalTok{(mask\_ras), }\ConstantTok{NA}\NormalTok{, }\DecValTok{1}\NormalTok{)}
\end{Highlighting}
\end{Shaded}

\section{Fossil matching}\label{fossil-matching}

The example fossil dataset \texttt{data(ex\_foss)} contains spatial
point locations and associated radiocarbon age estimates. The raw object
is stored as a SpatVector and needs to be converted into a data table
that can be passed to \texttt{pair\_obs}.

The \texttt{terra::unwrap} call ensures that all attributes are pulled
into the R session, including the age and confidence interval fields.

To prepare the data, the spatial coordinates are extracted from the
SpatVector, unique integer IDs are assigned, and the original date
estimate and its upper and lower confidence bounds are retained. The
resulting data table contains one row per observation, with explicit
latitude, longitude, and temporal uncertainty.

\begin{Shaded}
\begin{Highlighting}[]
\FunctionTok{data}\NormalTok{(ex\_foss)}
\NormalTok{ex\_foss }\OtherTok{\textless{}{-}}\NormalTok{ terra}\SpecialCharTok{::}\FunctionTok{unwrap}\NormalTok{(ex\_foss)}

\NormalTok{ex\_foss }\OtherTok{\textless{}{-}}\NormalTok{ data.table}\SpecialCharTok{::}\FunctionTok{data.table}\NormalTok{(}
  \AttributeTok{ID =} \DecValTok{1}\SpecialCharTok{:}\FunctionTok{nrow}\NormalTok{(ex\_foss),}
  \AttributeTok{Lat =} \FunctionTok{crds}\NormalTok{(ex\_foss)[,}\DecValTok{2}\NormalTok{],}
  \AttributeTok{Lon =} \FunctionTok{crds}\NormalTok{(ex\_foss)[,}\DecValTok{1}\NormalTok{],}
  \AttributeTok{Age =}\NormalTok{ ex\_foss}\SpecialCharTok{$}\NormalTok{date\_estimate,}
  \AttributeTok{AgeMin =}\NormalTok{ ex\_foss}\SpecialCharTok{$}\NormalTok{ci\_upper,}
  \AttributeTok{AgeMax =}\NormalTok{ ex\_foss}\SpecialCharTok{$}\NormalTok{ci\_lower}
\NormalTok{)}
\end{Highlighting}
\end{Shaded}

\begin{verbatim}
#>       ID     Lat     Lon    Age AgeMin AgeMax
#>    <int>   <num>   <num>  <num>  <num>  <num>
#> 1:     1 -11.575 132.725 -20900 -20859 -20941
#> 2:     2 -12.575 134.275 -20473 -20441 -20505
#> 3:     3 -12.225 133.875 -18041 -17984 -18097
#> 4:     4 -12.525 130.675 -17595 -17551 -17639
#> 5:     5 -12.725 130.475 -17259 -17233 -17285
#> 6:     6 -12.075 133.375 -16878 -16847 -16908
\end{verbatim}

The expected input structure for \texttt{pair\_obs} is a
\texttt{data.table} containing:

\begin{enumerate}
\def\labelenumi{\arabic{enumi}.}
\item
  Lat and Lon in geographic coordinates (decimal degrees)
\item
  Age, the central age estimate
\item
  AgeMin and AgeMax, representing the youngest and oldest bounds of the
  age range
\item
  any optional additional columns retained for downstream use
\end{enumerate}

\texttt{pair\_obs} uses the temporal range (AgeMin--AgeMax) to determine
which climate layers should be included during matching.

\subsection{\texorpdfstring{Environmental matching with
\texttt{pair\_obs}}{Environmental matching with pair\_obs}}\label{environmental-matching-with-pair_obs}

\texttt{pair\_obs} extracts climate values from time-indexed raster
stacks and associates them with each fossil record, accounting for
temporal uncertainty, spatial neighbourhoods, and optional
distance-based filtering. The function is designed specifically for
datasets like TraCE--Sahul, where each raster layer corresponds to a
single year (or summary period) and a long continuous time series is
available.

The matching workflow includes the following steps:

\begin{itemize}
\item
  Identification of candidate climate layers based on fossil ages and
  the specified temporal window (window). For example, a 30-year window
  uses the 30 most recent layers preceding the most-recent age estimate
  (i.e.~right aligned).
\item
  Spatial filtering using a mask (mask\_layer), which ensures that only
  land pixels or valid climate surfaces are used. Records occurring on
  land/sea can be filtered post-matching depending on what is required.
\item
  Extraction of nearby cell values using the requested neighbourhood
  size (\texttt{neigh}). Values may be taken from the closest cell, or
  averaged across all neighbouring cells if a summary statistic is
  supplied.
\item
  Optional distance-based exclusion using \texttt{dist\_cut}, which
  removes records where the spatial offset between the fossil location
  and the raster grid exceeds a specified threshold.
\item
  Precision control through the prec argument, which rounds output
  values.
\item
  Parallelisation using cores for improved performance on large
  datasets. All rasters are wrapped and unwrapped internally to ensure
  compatibility with parallel workers. This is currently disabled.
\end{itemize}

The function returns a data table containing the original fossil
metadata, the matched climate values, the temporal layer selected for
each match, and diagnostic information about mask status, neighbourhood
selection, and any records that could not be matched (e.g.~ages older
than the available climate layers).

In the following example, climate values are drawn from BIOCLIM rasters
(\texttt{ras\_list\ =\ bioclims}). A 30-year right-aligned temporal
window is used (\texttt{window\ =\ 30}), eight-cell neighbourhoods are
searched (\texttt{neigh\ =\ 8}), and values are averaged
(\texttt{summ\_stat\ =\ "mean"}) when multiple neighbouring cells are
available. Records more than 50 km from the raster grid are excluded
(\texttt{dist\_cut\ =\ 50}).

The final output contains matched BIOCLIM variables for all fossils
whose ages fall within the time span of the dataset.

\begin{Shaded}
\begin{Highlighting}[]
\NormalTok{foss\_matched\_bioclim }\OtherTok{\textless{}{-}} \FunctionTok{pair\_obs}\NormalTok{(}
  \AttributeTok{data =}\NormalTok{ ex\_foss,}
  \AttributeTok{ras\_list =}\NormalTok{ bioclims,}
  \AttributeTok{mask\_layer =}\NormalTok{ mask\_ras,}
  \AttributeTok{ras\_time =}\NormalTok{ terra}\SpecialCharTok{::}\FunctionTok{time}\NormalTok{(bioclims}\SpecialCharTok{$}\NormalTok{bio01),}
  \AttributeTok{window =} \DecValTok{30}\NormalTok{, }\CommentTok{\# 30{-}year window}
  \AttributeTok{neigh =} \DecValTok{8}\NormalTok{, }\CommentTok{\# 8{-}neighbours (rook)}
  \AttributeTok{prec =} \DecValTok{3}\NormalTok{, }\CommentTok{\# 3 decimal places}
  \AttributeTok{summ\_stat =} \StringTok{"mean"}\NormalTok{,}
  \AttributeTok{dist\_cut =} \DecValTok{50}\NormalTok{, }\CommentTok{\# km}
  \CommentTok{\# parallel is disabled and will be set to 1 if any other value is given}
  \AttributeTok{cores =} \DecValTok{1}\DataTypeTok{L}\NormalTok{, }
  \AttributeTok{buff\_width =} \ConstantTok{NULL}
\NormalTok{)}
\CommentTok{\#\textgreater{} Lon/Lat points will be snapped to nearest points on rasters. Cutoff distance = 50 km.}
\CommentTok{\#\textgreater{} Extracting data from rasters sequentially}
\CommentTok{\#\textgreater{} Warning in pair\_obs(data = ex\_foss, ras\_list = bioclims, mask\_layer = mask\_ras, : }
\CommentTok{\#\textgreater{} There were 2 samples removed (no temporal match, \textgreater{} cutoff distance, or extraction error).}
\CommentTok{\#\textgreater{} Warning in pair\_obs(data = ex\_foss, ras\_list = bioclims, mask\_layer = mask\_ras, : Samples removed: 1, 2}
\end{Highlighting}
\end{Shaded}

\section{Species distribution
modelling}\label{species-distribution-modelling}

Following data aggregation and record matching, we can create a simple
species distribution model to examine how habitat suitability (predicted
occurence) has changed through time.

Here we use downsampled random forests after extracting some background
values through time.

The down-sampling approach for random forests ensures that each tree in
the Random Forest is trained using an equal number of presence and
background samples. This is controlled through the \texttt{sampsize}
argument, which accepts a named vector specifying how many observations
to draw from each class during bootstrap sampling. Although each tree is
fitted with balanced class sizes, the final forest still incorporates
the full set of background records because different subsets are
selected for each tree across the ensemble.

\subsection{Background point
generation}\label{background-point-generation}

Here we generate 12,600 background points (50 samples x 252 timesteps)

\begin{Shaded}
\begin{Highlighting}[]
\CommentTok{\# which timesteps from the TraCE{-}Sahul data are included in our fossil matched data}
\NormalTok{tidx }\OtherTok{\textless{}{-}} \FunctionTok{which}\NormalTok{(terra}\SpecialCharTok{::}\FunctionTok{time}\NormalTok{(bioclims}\SpecialCharTok{$}\NormalTok{bio01) }\SpecialCharTok{\%in\%}\NormalTok{ foss\_matched\_bioclim}\SpecialCharTok{$}\NormalTok{Year)}

\CommentTok{\# extract bg points.}
\DocumentationTok{\#\# seed is set for reproducability}
\NormalTok{\{}\FunctionTok{set.seed}\NormalTok{(}\DecValTok{84564}\NormalTok{)}
\NormalTok{bg\_points }\OtherTok{\textless{}{-}} \FunctionTok{lapply}\NormalTok{(tidx, }\ControlFlowTok{function}\NormalTok{(t) \{}
\NormalTok{  nsamps }\OtherTok{\textless{}{-}} \DecValTok{50}
\NormalTok{  preds }\OtherTok{\textless{}{-}} \FunctionTok{rast}\NormalTok{(}\FunctionTok{lapply}\NormalTok{(}\FunctionTok{seq\_along}\NormalTok{(bioclims), }\ControlFlowTok{function}\NormalTok{(i,...) bioclims[[i]][[t]]))}
  \FunctionTok{names}\NormalTok{(preds) }\OtherTok{\textless{}{-}} \FunctionTok{sapply}\NormalTok{(}\FunctionTok{strsplit}\NormalTok{(}\FunctionTok{names}\NormalTok{(preds), }\StringTok{"\_"}\NormalTok{), }\StringTok{"["}\NormalTok{, }\DecValTok{1}\NormalTok{)}
\NormalTok{  data.table}\SpecialCharTok{::}\FunctionTok{setDT}\NormalTok{(terra}\SpecialCharTok{::}\FunctionTok{spatSample}\NormalTok{(preds, nsamps, }\AttributeTok{na.rm =} \ConstantTok{TRUE}\NormalTok{))}
\NormalTok{\})\}}

\CommentTok{\# collate the sample and round}
\NormalTok{bg\_points }\OtherTok{\textless{}{-}}\NormalTok{ data.table}\SpecialCharTok{::}\FunctionTok{rbindlist}\NormalTok{(bg\_points)}
\NormalTok{cols }\OtherTok{\textless{}{-}} \FunctionTok{names}\NormalTok{(bg\_points)}
\NormalTok{bg\_points[,(cols) }\SpecialCharTok{:=} \FunctionTok{round}\NormalTok{(.SD, }\DecValTok{3}\NormalTok{), .SDcols }\OtherTok{=}\NormalTok{ cols]}
\end{Highlighting}
\end{Shaded}

A quick looks shows what we generated for the first 5 variables

\begin{verbatim}
#>     bio01   bio04  bio05  bio06  bio07
#>     <num>   <num>  <num>  <num>  <num>
#> 1: 24.970  64.277 27.839 22.126  5.713
#> 2: 25.362  59.277 27.267 23.515  3.753
#> 3: 25.340  59.161 27.158 23.661  3.496
#> 4: 25.333  60.906 27.235 23.503  3.733
#> 5: 25.353 196.556 34.818 16.446 18.372
#> 6: 25.466  68.488 28.073 22.760  5.313
\end{verbatim}

\subsection{Random Forest modelling}\label{random-forest-modelling}

Here we generate the training data and build the random forest.

Columns not needed are excluded, and the presence and absence are set as
factors so we can do classification modelling.

\begin{Shaded}
\begin{Highlighting}[]
\NormalTok{training }\OtherTok{\textless{}{-}}\NormalTok{ data.table}\SpecialCharTok{::}\FunctionTok{rbindlist}\NormalTok{(}\FunctionTok{list}\NormalTok{(}
\NormalTok{  data.table}\SpecialCharTok{::}\FunctionTok{copy}\NormalTok{(foss\_matched\_bioclim)[, }\SpecialCharTok{{-}}\FunctionTok{c}\NormalTok{(}\StringTok{"ID"}\NormalTok{, }\StringTok{"Lon"}\NormalTok{, }\StringTok{"Lat"}\NormalTok{, }\StringTok{"Year"}\NormalTok{, }\StringTok{"LandSea"}\NormalTok{)][, }\AttributeTok{Occ :=} \StringTok{"1"}\NormalTok{],}
\NormalTok{  data.table}\SpecialCharTok{::}\FunctionTok{copy}\NormalTok{(bg\_points)[, }\AttributeTok{Occ :=} \StringTok{"0"}\NormalTok{]}
\NormalTok{))}

\NormalTok{training[, Occ }\SpecialCharTok{:=} \FunctionTok{factor}\NormalTok{(Occ, }\AttributeTok{levels =} \FunctionTok{c}\NormalTok{(}\StringTok{"0"}\NormalTok{,}\StringTok{"1"}\NormalTok{))]}

\NormalTok{prNum }\OtherTok{\textless{}{-}} \FunctionTok{as.numeric}\NormalTok{(}\FunctionTok{table}\NormalTok{(training}\SpecialCharTok{$}\NormalTok{Occ)[}\StringTok{"1"}\NormalTok{])}
\NormalTok{spsize }\OtherTok{\textless{}{-}} \FunctionTok{c}\NormalTok{(}\StringTok{"0"} \OtherTok{=}\NormalTok{ prNum, }\StringTok{"1"} \OtherTok{=}\NormalTok{ prNum)}

\NormalTok{rf\_dws }\OtherTok{\textless{}{-}} \FunctionTok{randomForest}\NormalTok{(}
\NormalTok{  Occ }\SpecialCharTok{\textasciitilde{}}\NormalTok{ .,}
  \AttributeTok{mtry =} \DecValTok{5}\NormalTok{,}
  \AttributeTok{data =}\NormalTok{ training,}
  \AttributeTok{ntree =} \DecValTok{1000}\NormalTok{,}
  \AttributeTok{importance =} \ConstantTok{TRUE}\NormalTok{,}
  \AttributeTok{sampsize =}\NormalTok{ spsize,}
  \AttributeTok{replace =} \ConstantTok{TRUE}\NormalTok{)}
\end{Highlighting}
\end{Shaded}

We can see that the random forest model was moderatly accurate using the
misclassification matrix

\begin{verbatim}
#> 
#> Call:
#>  randomForest(formula = Occ ~ ., data = training, mtry = 5, ntree = 1000,      importance = TRUE, sampsize = spsize, replace = TRUE) 
#>                Type of random forest: classification
#>                      Number of trees: 1000
#> No. of variables tried at each split: 5
#> 
#>         OOB estimate of  error rate: 15.8%
#> Confusion matrix:
#>       0    1 class.error
#> 0 10609 1991   0.1580159
#> 1    48  259   0.1563518
\end{verbatim}

\subsection{Raster prediction example}\label{raster-prediction-example}

Finally, we can predict from the RF model to generate layers of
prediction probability (or habitat suitability).

Here we will extract the rasters and predict at a number of timesteps so
we can see how probability of occurrence changes through time.

\begin{Shaded}
\begin{Highlighting}[]
\CommentTok{\# time steps to extract at}
\NormalTok{ts }\OtherTok{\textless{}{-}} \FunctionTok{c}\NormalTok{(}\DecValTok{1}\NormalTok{, }\DecValTok{65}\NormalTok{, }\DecValTok{422}\NormalTok{, }\DecValTok{507}\NormalTok{, }\DecValTok{730}\NormalTok{)}

\CommentTok{\# HS predictions}
\NormalTok{preds\_stack }\OtherTok{\textless{}{-}} \FunctionTok{rast}\NormalTok{(}\FunctionTok{lapply}\NormalTok{(ts, }\ControlFlowTok{function}\NormalTok{(t) \{}
\NormalTok{  r }\OtherTok{\textless{}{-}} \FunctionTok{rast}\NormalTok{(}\FunctionTok{lapply}\NormalTok{(}\FunctionTok{seq\_along}\NormalTok{(bioclims), }\ControlFlowTok{function}\NormalTok{(i) \{}
\NormalTok{    bioclims[[i]][[t]]}
\NormalTok{  \}))}
  \CommentTok{\# Give the correct names to the layers in each stack}
  \FunctionTok{names}\NormalTok{(r) }\OtherTok{\textless{}{-}} \FunctionTok{sapply}\NormalTok{(}\FunctionTok{strsplit}\NormalTok{(}\FunctionTok{names}\NormalTok{(r), }\StringTok{"\_"}\NormalTok{), }\StringTok{"["}\NormalTok{, }\DecValTok{1}\NormalTok{)}
  \CommentTok{\# predict prob of occurence}
\NormalTok{  pr }\OtherTok{\textless{}{-}} \FunctionTok{predict}\NormalTok{(r, rf\_dws, }\AttributeTok{type =} \StringTok{"prob"}\NormalTok{, }\AttributeTok{index =} \DecValTok{2}\NormalTok{)}
  \FunctionTok{names}\NormalTok{(pr) }\OtherTok{\textless{}{-}}\NormalTok{ terra}\SpecialCharTok{::}\FunctionTok{time}\NormalTok{(bioclims}\SpecialCharTok{$}\NormalTok{bio01)[t]}
  \FunctionTok{return}\NormalTok{(pr)}
\NormalTok{\}))}

\NormalTok{preds\_stack}
\CommentTok{\#\textgreater{} class       : SpatRaster }
\CommentTok{\#\textgreater{} size        : 100, 200, 5  (nrow, ncol, nlyr)}
\CommentTok{\#\textgreater{} resolution  : 0.05, 0.05  (x, y)}
\CommentTok{\#\textgreater{} extent      : 125, 135, {-}13, {-}8  (xmin, xmax, ymin, ymax)}
\CommentTok{\#\textgreater{} coord. ref. : lon/lat WGS 84 (EPSG:4326) }
\CommentTok{\#\textgreater{} source(s)   : memory}
\CommentTok{\#\textgreater{} names       : {-}20015, {-}18095, {-}7385, {-}4835,  1869 }
\CommentTok{\#\textgreater{} min values  :  0.006,  0.000, 0.000, 0.001, 0.002 }
\CommentTok{\#\textgreater{} max values  :  0.801,  0.864, 0.931, 0.897, 0.632}
\end{Highlighting}
\end{Shaded}

Now we can get the \emph{true} suitability and add it to our predicted
suitability for comparison.

\begin{Shaded}
\begin{Highlighting}[]
\CommentTok{\# get the true\_suit raster}
\FunctionTok{data}\NormalTok{(true\_suit)}
\NormalTok{true\_suit }\OtherTok{\textless{}{-}} \FunctionTok{unwrap}\NormalTok{(true\_suit)}
\FunctionTok{names}\NormalTok{(true\_suit) }\OtherTok{\textless{}{-}} \StringTok{"true suitability"}

\CommentTok{\# add the true\_suit to the preds\_stack}
\NormalTok{preds\_stack }\OtherTok{\textless{}{-}} \FunctionTok{c}\NormalTok{(preds\_stack, true\_suit)}
\NormalTok{preds\_stack}
\CommentTok{\#\textgreater{} class       : SpatRaster }
\CommentTok{\#\textgreater{} size        : 100, 200, 6  (nrow, ncol, nlyr)}
\CommentTok{\#\textgreater{} resolution  : 0.05, 0.05  (x, y)}
\CommentTok{\#\textgreater{} extent      : 125, 135, {-}13, {-}8  (xmin, xmax, ymin, ymax)}
\CommentTok{\#\textgreater{} coord. ref. : lon/lat WGS 84 (EPSG:4326) }
\CommentTok{\#\textgreater{} source(s)   : memory}
\CommentTok{\#\textgreater{} names       : {-}20015, {-}18095, {-}7385, {-}4835,  1869, true \textasciitilde{}ility }
\CommentTok{\#\textgreater{} min values  :  0.006,  0.000, 0.000, 0.001, 0.002,           0 }
\CommentTok{\#\textgreater{} max values  :  0.801,  0.864, 0.931, 0.897, 0.632,           1}
\end{Highlighting}
\end{Shaded}

Now we can plot the raster stack of predictions and compare it to the
\emph{true} habitat suitability that was used to generate the virtual
species occurrences in the \texttt{ex\_foss} records.

\begin{center}\includegraphics[width=1\linewidth]{TraCESahulMisc_workflow_files/figure-latex/plotPreds-1} \end{center}

We can see that there are changes in habitat suitability through time.
It's also obvious that the Random Forest SDM is not great at reproducing
the \emph{true} suitability of the virtual species, however this is to
be expected as the virtual species used a different modelling approach
and variables!

\end{document}
